\section{Conclusion}
\label{sec:conclusion}

% Remaster, 5 October 2026.  Shortened; numbers stay in the Findings.

We asked where the accuracy of LLM-agent AutoML comes from and measured
three candidate sources separately: the harness, the instruction file and
the backbone. One open-weight 31B model ran under three harnesses on eight
real-world tabular tasks with temporal splits, and eighteen
hyperparameter-tuned reference methods served as the yardstick.

With the instruction file removed, a fixed-stage workflow, a six-agent
system and a minimal shell agent lie within $\spanArch$ of each other on the
normalised scale, no more than the variation between attempts of one
harness, and no pairwise difference is significant. The harness does affect
reliability, wall-clock time and the set of tasks a system can attempt.

The instruction file accounts for the accuracy gain, and its effect depends
on the harness and on the starting point. Used from scratch inside the
harness it was written for, it improves all eight tasks and brings the
system on par with tuned LightGBM. Appended to another harness, it ends half
or more of the trials without a submission; where a working script already
exists, it leaves whole tasks without one. An instruction file is therefore
part of the system under test and must be evaluated in each harness and from
each starting point.

No agent searches over hyperparameters, and this separates every
configuration from the strongest published method. A leaderboard built
around one harness measures the instructions that ship with it as well as
the model. Two experiments follow from the traces and remain to be run:
adding hyperparameter search to each harness under the same budget, and
disabling the repair loop of AutoDS-Tools to measure the cost of the
instruction file without it.

\section*{Data availability}

\begin{sloppypar}
All tasks, the instruction file in every variant used, and every execution
trace behind these numbers are public. The benchmark adaptations and the instruction files are
on HuggingFace as \texttt{danil-e/harbor-datasets-tabred} and
\texttt{danil-e/harbor-datasets-mlab}; the tabular adapter is
\texttt{AaLexUser/harbor-tabred-adapter}. Trajectories, submissions, scorer
outputs and telemetry of the earlier jobs are on the Harbor hub under job
identifiers \texttt{7593cc7b}, \texttt{1a42d1fe}, \texttt{3113aba7} and
\texttt{56768b37}, and the FEDOT.LLM TabReD row under \texttt{77c2d2f8},
\texttt{374be194}, \texttt{b9de44ad} and \texttt{7768c071}, where the same
tasks and scorer appear under the dataset name \texttt{tabred-fedot-schema},
the copy built on the FEDOT image; the repeats are
released in the data package with the
consolidated metric table, the per-task tables behind every summary, the
learner extraction of Section~\ref{sec:spontaneous}, the telemetry of
Table~\ref{tab:cost}, the byte lengths and hashes of every instruction-file variant, and the scripts
that produce every table and figure.
\end{sloppypar}
\todo{repository DOI once deposited}

\section*{CRediT authorship contribution statement}

\textbf{Stanislav Chumakov:} Conceptualization, Methodology, Investigation,
Formal analysis, Software, Validation, Visualization, Writing -- original
draft, Writing -- review \& editing.
\textbf{Danil Ezhov:} Investigation (AutoDS-Tools experiments), Software,
Data curation.
\textbf{Aleksei Lapin:} Investigation (Terminus-2 experiments), Software.
\textbf{Pavel Marian:} Investigation (FEDOT.LLM experiments), Software,
Resources (experiment tracking in Harbor).
\textbf{Kirill Anpilov:} Investigation (scientific case studies), Data
curation (case-study selection and task construction).
\textbf{Nikolay O. Nikitin:} Conceptualization, Supervision, Project
administration, Funding acquisition, Writing -- review \& editing.

\section*{Declaration of competing interest}

The authors declare that they have no known competing financial interests or
personal relationships that could have appeared to influence the work reported
in this paper.

\section*{Declaration of generative AI and AI-assisted technologies in the manuscript preparation process}

During the preparation of this work the authors used Claude Code (Anthropic)
to assist with writing and debugging the scripts that run the experiments and
produce the tables and figures, and to edit and format the manuscript text.
All tables and figures are generated by those scripts from the deposited
results; no image was created or altered by a generative model. After using this tool, the authors
reviewed and edited the content as needed and take full responsibility for the
content of the published article.

\section*{Acknowledgements}
\todo{funding, compute, and people to thank}
